% Generated draft: requires mathematical and source review.
\documentclass[11pt]{article}
\usepackage[margin=25mm]{geometry}
\usepackage{fontspec}
\setmainfont{Latin Modern Roman}
\usepackage{amsmath,amssymb,mathtools,unicode-math}
\setmathfont{Latin Modern Math}
\usepackage{enumitem,needspace,xurl,hyperref}
\setlength{\parindent}{0pt}
\setlength{\parskip}{5pt}
\setlength{\emergencystretch}{3em}
\newcommand{\N}{\mathbb{N}}
\newcommand{\Z}{\mathbb{Z}}
\newcommand{\Q}{\mathbb{Q}}
\newcommand{\R}{\mathbb{R}}
\newcommand{\C}{\mathbb{C}}
\newcommand{\F}{\mathbb{F}}
\newcommand{\A}{\mathbb{A}}
\newcommand{\PP}{\mathbb{P}}
\newcommand{\E}{\mathbb{E}}
\newcommand{\eps}{\varepsilon}
\newcommand{\dd}{\,\mathrm{d}}
\begin{document}
\section*{UnsolvedMath ID: 30002250}
\textbf{Source problem code:} catalogue entry
\subsection{Definitions and mathematical statement}
Let \(n\geq 1\) be an integer. The source considers convex bodies \(K,E\in\mathcal K^n\) in \(\R^n\) and a nonnegative measure \(\mu\) on \([0,\infty)\). It does not specify whether its convention for convex bodies requires nonempty interior.

Under the usual relative mixed-volume normalization, the coefficients \(W_i(K;E)\) are characterized by
\[
\operatorname{vol}_n(K+tE)=\sum_{i=0}^{n}\binom{n}{i}W_i(K;E)t^i,\qquad t\geq0,
\]
where \(K+tE=\{x+ty:x\in K,\ y\in E\}\). This normalization is an interpretive convention; the supplied record does not define \(W_i\).

The displayed expression in the source is
\[
g_{K;E}^{\mu}(z)=\sum_{i=0}^{n}\binom{n}{i}W_i(K;E)m_i(\mu)z^i.
\]
The quantities \(m_i(\mu)\) are not defined in the supplied record. If ordinary moments are intended, their definition would be
\[
m_i(\mu)=\int_{[0,\infty)}t^i\,\mathrm d\mu(t),\qquad 0\leq i\leq n,
\]
with \(t^0=1\), including at \(t=0\). This interpretation requires confirmation. Under it, finiteness of all these moments is a sufficient additional hypothesis for the displayed coefficients to be finite. Nonnegativity of \(\mu\) alone does not provide this hypothesis.

Whenever the coefficients are well-defined and finite and \(g_{K;E}^{\mu}\not\equiv0\), let \(\mathcal Z_{K,E}^{\mu}\) be the multiset of its complex roots, counted with algebraic multiplicity. Its degree may be smaller than \(n\); a nonzero constant polynomial has an empty root multiset. Under the ordinary-moment interpretation, the allowed zero measure produces the zero polynomial, whose roots cannot be treated as a finite multiset counted with algebraic multiplicity.

\textbf{Recorded research aim, with unresolved formulation.} Determine structural properties shared by \(\mathcal Z_{K,E}^{\mu}\) independently of the admissible choice of \(\mu\). The source does not identify the properties to be characterized, the admissible measure class, or whether the intended uniformity fixes \(n,K,E\) or also varies them. Consequently, this is a research aim rather than a specified proposition with complete quantifiers.

\textbf{Source qualification.} The supplied statement-verification note says that the official contribution introduces this research aim and then states properties it proves. It identifies the extraction as an answered research aim, not a separately posed open problem, and marks the statement as unrecoverable. No separate clean statement is supplied. This draft therefore records the aim and its missing hypotheses without promoting it to a verified open conjecture.
\subsection{Short English statement}
The record asks which features of generalized Wills-polynomial roots persist as the measure varies. The moment convention, admissible measures, and precise target properties are missing. Supplied verification identifies the passage as an answered research aim rather than a separate open problem; the supplied partially solved label is dated curation, not an independently verified conclusion.
\subsection{Sources}
\begin{itemize}
\item Supplied local UnsolvedMath registry record, permanent ID \(30002250\), title \emph{Generalized Wills-Polynomial Results Report}; source problem code \texttt{catalogue entry}. The record belongs to the Oberwolfach Reports Open Problems collection. No individual proposer is supplied; the recorded proposed year is 2012.
\item Supplied report citation: \emph{Convex Geometry and its Applications (2013)}, Oberwolfach Reports, pp.~3473--3548, DOI \texttt{10.4171/owr/2012/59}: \url{https://doi.org/10.4171/owr/2012/59}. The background describes Oberwolfach Reports as documenting mathematical workshops and their open-problem discussions. The citation's 2013 designation and its 2012 DOI component are retained as supplied.
\item Supplied literature source, labeled ``Oberwolfach report'': \url{https://ems.press/content/serial-article-files/46430}.
\item Supplied literature source, labeled ``Wills polynomial roots theorem'': \url{https://doi.org/10.4171/RMI/842}. No author names or more detailed bibliographic information for this item are supplied.
\item Literature assessment checked at \texttt{2026-08-21}: the supplied curation labels the record partially solved and reports substantial measure-independent structure, including convex cone containment and monotonicity properties, while stating that the broad request for all common properties is not a complete classification. This is supplied dated curation derived from the OWR material, not independently verified here; its time-sensitive status and cited claims require independent verification.
\item Supplied statement-verification note: the official contribution introduces interest in measure-independent root properties and then states properties it proves; the extraction is an answered research aim, not a separately posed open problem. The record marks the statement as unrecoverable, promotes no separate clean statement, and retains the original publicly visible record for provenance. No separate verification date is supplied.
\end{itemize}
\end{document}
